\documentclass[11pt,letterpaper]{article}
\usepackage[margin=1in,headheight=14pt]{geometry}
\usepackage[T1]{fontenc}
\usepackage{amsmath,amssymb,amsthm}
\usepackage{newtxtext,newtxmath}
\usepackage{microtype}
\usepackage{booktabs,array,longtable}
\usepackage{mathtools}
\usepackage{xcolor}
\usepackage{enumitem}
\usepackage{listings}
\usepackage{fancyhdr}
\usepackage[hyphens]{url}
\usepackage{hyperref}
\hypersetup{colorlinks=true,linkcolor=blue!45!black,citecolor=blue!45!black,
 urlcolor=blue!45!black,pdftitle={A Cayley Certificate for the Mixed Gaussian S4 Reduction},
 pdfauthor={ProveIt research continuation}}
\numberwithin{equation}{section}
\newtheorem{theorem}{Theorem}[section]
\newtheorem{lemma}[theorem]{Lemma}
\newtheorem{proposition}[theorem]{Proposition}
\newtheorem{corollary}[theorem]{Corollary}
\newtheorem{conjecture}[theorem]{Conjecture}
\theoremstyle{definition}
\newtheorem{definition}[theorem]{Definition}
\newtheorem{remark}[theorem]{Remark}
\DeclareMathOperator{\Li}{Li}
\DeclareMathOperator{\Imag}{Im}
\DeclareMathOperator{\Real}{Re}
\DeclareMathOperator{\TV}{TV}
\DeclareMathOperator{\Odd}{Odd}
\newcommand{\Q}{\mathbb Q}
\newcommand{\Z}{\mathbb Z}
\newcommand{\C}{\mathbb C}
\newcommand{\R}{\mathbb R}
\newcommand{\dd}{\,\mathrm d}
\newcommand{\ii}{\mathrm i}
\newcommand{\sh}{\mathbin{\sha}}
% A font-independent shuffle symbol.
\providecommand{\sha}{\mathbin{\sqcup\!\sqcup}}
\newcommand{\eps}{\varepsilon}
\newcommand{\cay}{\mathcal D}
\newcommand{\word}{\operatorname{word}}
\newcommand{\ind}{\operatorname{ind}}
\newcommand{\norm}[1]{\left\lVert#1\right\rVert}
\setlist{topsep=4pt,itemsep=3pt,parsep=0pt}
\lstset{basicstyle=\ttfamily\footnotesize,columns=fullflexible,keepspaces=true,
 breaklines=true,frame=single,rulecolor=\color{black!20},showstringspaces=false,
 aboveskip=10pt,belowskip=10pt}
\pagestyle{fancy}
\fancyhf{}
\fancyhead[L]{\small Cayley certificates and Gaussian polylogarithms}
\fancyhead[R]{\small ProveIt research continuation}
\fancyfoot[C]{\thepage}
\setlength{\parskip}{3pt}
\setlength{\emergencystretch}{2em}

\title{\textbf{A Cayley Certificate for the\\Mixed Gaussian $S_4$ Reduction}\\[7pt]
\large Exact level-four identities, signed-moment certification,\\
and a weight-seven conjecture}
\author{Prepared for Vladimir Reshetnikov\\
\small A research continuation for the ProveIt project}
\date{9 October 2026}

\begin{document}
\maketitle
\begin{abstract}
We prove the mixed Gaussian identity left conjectural in the inspected
ProveIt polylogarithm manuscript,
\[
\sum_{n\ge0}\frac{(-1)^nH_n}{(2n+1)^4}
 =\frac{58\Imag\Li_{4,1}(i,1)+24\Imag\Li_{3,2}(i,1)}7
  +\frac{19\pi^5}{3584}-2\beta(4)\log2.
\]
The proof is computer-assisted but entirely exact: a supplied rational
certificate writes a specified weight-five word polynomial as a combination
of 852 double-shuffle relations and one relation induced by the Cayley
involution $t\mapsto(1-t)/(1+t)$. We establish the analytic validity of every
relation type, including the 123 rows with one initially divergent factor.
Two separately implemented verifiers check the finite word identity without
floating-point arithmetic or a linear solver. The result also gives an
explicit reduction of a mixed-color antisymmetric double.

For reproducible experimentation, we prove a signed-Hausdorff-moment bound
for the more general coefficients $H_{rk}^{(b)}/(dk+1)^a$, with an explicit
Gamma-density improvement depending on $a$. An exact rational Euler
transform then certifies errors geometrically. It supports a new proposed
weight-seven identity for $S_6$ with a residual enclosure of width below
$3\cdot10^{-299}$, while rigorously excluding a misleading weight-nine
near-relation. We identify the precise manuscript status changes warranted
by the proof, distinguish finite relation spaces from arithmetic
independence, and propose further research. The Cayley symmetry,
double-shuffle framework, and Euler transform are established tools;
no worldwide-priority claim is made for them or for the special-value
formulas. The new contribution to the inspected project is the explicit
proof certificate, its consequences, and the documented continuation.
\end{abstract}

\vfill
\noindent\textbf{Proof status.} The $S_4$ identity and the analytic error bounds
are proved. The $S_6$ identity remains a conjecture. The displayed $S_8$
near-relation is proved false. There is no proof-assistant formalization.
\newpage
\begingroup
\small
\setlength{\parskip}{0pt}
\tableofcontents
\endgroup
\newpage

\section{Research target, provenance, and main results}
\label{sec:intro}

\subsection{The precise starting point}
The project source is the collective manuscript \emph{Polylogarithms and
their Arithmetic Bridges} in
\path{Analysis/Polylogarithms/docs/manuscript} of the ProveIt repository
\cite{ProveIt}. The working snapshot is commit
\begin{center}
\small\texttt{a2a4cf58c49c745058c40e4a6748d472a3420f18}.
\end{center}
The audit here is targeted, not a claim to have independently verified the
entire book. The relevant starting material is the mixed relation in
\path{chapters/04-shuffle-parity.tex}, with source labels
\texttt{gaussian:conj:S4} and \texttt{gaussian:eq:S4-short}. That source already
proves the ordinary weight-five Gaussian shuffle reductions and explicitly
leaves the mixed $S_4$ identity conjectural. We retain its notation:
\begin{equation}
 H_n^{(b)}=\sum_{j=1}^n\frac1{j^b},\quad H_n=H_n^{(1)},\quad H_0^{(b)}=0,
 \qquad S_p=\sum_{n=0}^{\infty}\frac{(-1)^nH_n}{(2n+1)^p},
 \label{eq:Sp}
\end{equation}
\begin{equation}
 g_{a,b}=\Imag\Li_{a,b}(i,1),\qquad
 \beta(s)=\sum_{n=0}^{\infty}\frac{(-1)^n}{(2n+1)^s},\qquad L=\log2.
 \label{eq:notation}
\end{equation}
All indices in the special-value identities below are positive integers.
The sums with $p>1$ are absolutely convergent; the analytic arguments also
cover several conditionally convergent index-one values needed by the word
relations.

\begin{theorem}[The mixed Gaussian $S_4$ identity]\label{thm:S4}
In the strict nested-sum convention defined in Section~\ref{sec:words},
\begin{equation}
 \boxed{\displaystyle
 S_4=\frac{58g_{4,1}+24g_{3,2}}7+\frac{19\pi^5}{3584}
       -2\beta(4)L.}
 \label{eq:S4}
\end{equation}
Equivalently,
\begin{equation}
 \Imag\Li_{4,1}(i,-1)
 =\frac{51g_{4,1}+24g_{3,2}}7+\frac{19\pi^5}{3584}-2\beta(4)L.
 \label{eq:mixed41}
\end{equation}
\end{theorem}
The proof occupies Sections~\ref{sec:regularization}--\ref{sec:certificate}.
Its finite arithmetic component is supplied in full in the accompanying
\path{data/S4_certificate.json}; it is not an integer-relation guess.

\subsection{What is added, and what is not claimed}
The important extension is the use of a symmetry relation that changes
intermediate depth. The single seed is the word for $\Li_{4,1}(i,1)$, but
its Cayley transform contains words of depth five. Thus the calculation is
not confined to an initially selected list of depth-two coordinates.
Ordinary shuffle and stuffle subsequently eliminate those intermediate
words. This explains the architecture of the proof without pretending
that the delivered 852-row double-shuffle part is minimal.

The paper also proves the exact mixed antisymmetric identity
\eqref{eq:antisymmetric}, an all-word Cayley formula with explicit endpoint
conditions, and a general signed-moment estimate. The latter produces
rational enclosures for both genuine and spurious numerical relations.
The $S_6$ formula in Conjecture~\ref{conj:S6} is a proposed new identity in
this continuation; its very small certified residual is not a proof that
it is exactly zero.

No dimension of a numerical period space is inferred. In particular, we
do not prove that $g_{4,1}$ and $g_{3,2}$ are independent, that depth two is
minimal, that the same basket suffices in every weight, or that the
Cayley row is outside the full ideal of all standard relations. The proof
uses only one explicitly defined collection of double-shuffle rows; it
is not an exhaustive comparison with every possible lifted, distribution,
or weight-one relation.

\subsection{Relationship to established methods}
Zhao's level-four work already uses the octahedral symmetry of the punctures
$\{0,\infty,1,-1,i,-i\}$ to obtain relations not detected by certain
standard-relation calculations \cite{Zhao2008,Zhao2010}. The present Cayley
map belongs to that established symmetry framework. Our task is the
specific manuscript identity, with explicit signs, a finite certificate,
and boundary-value justification. The much broader general theory of
regularized relations is not being rediscovered here.

Likewise, Euler acceleration is classical \cite{DLMF}. The analytic
contribution below is an explicit signed-measure bound for affine harmonic
coefficients, including an $a$-dependent refinement and a simple rational
implementation. The source manuscript already has signed-kernel and
Chebyshev-moment evaluation methods. We do not claim that the present
$2^{-N}$ bound improves the geometric rate of its Chebyshev method; its
advantage here is that every required special value can be enclosed using
rational harmonic numbers alone.

\section{Analytic conventions and the level-four word algebra}
\label{sec:words}

\subsection{Strict sums and cumulative colors}
For indices $s_j\ge1$ and colors $z_j$, write
\begin{equation}
 \Li_{s_1,\ldots,s_d}(z_1,\ldots,z_d)
 =\sum_{n_1>\cdots>n_d>0}
   \frac{z_1^{n_1}\cdots z_d^{n_d}}
        {n_1^{s_1}\cdots n_d^{s_d}}.
 \label{eq:MPL}
\end{equation}
At roots of unity the values used here are the ordinary convergent sums,
or equivalently their radial Abel limits. It is essential that the
\emph{series colors} in \eqref{eq:MPL} are relative colors. The corresponding
nonzero letters of an iterated integral are cumulative products.

Use the forms
\begin{equation}
 Z=\frac{\dd t}{t},\qquad e_j=\frac{i^j\dd t}{1-i^jt},\qquad j\in\Z/4\Z.
 \label{eq:forms}
\end{equation}
For a word $w=a_1\cdots a_n$, the convention is outer-letter first:
\begin{equation}
 I(aw;x)=\int_0^x\omega_a(t)I(w;t),\qquad I(\varnothing;x)=1,
 \qquad I(w)=I(w;1).
 \label{eq:iterated}
\end{equation}
A nonempty word is called \emph{admissible} if its last letter is not $Z$
and its first letter is not $e_0$. These conditions control the endpoints
at zero and one, respectively. Admissible iterated integrals converge
absolutely as ordered integrals, even when a corresponding outer-index-one
series converges only conditionally.

If $z_j=i^{c_j}$, then
\begin{equation}
 \word\bigl((s_1,c_1),\ldots,(s_d,c_d)\bigr)
 =Z^{s_1-1}e_{c_1}Z^{s_2-1}e_{c_1+c_2}\cdots
   Z^{s_d-1}e_{c_1+\cdots+c_d}.
 \label{eq:wordmap}
\end{equation}
Expanding the geometric kernels in the disk, or differentiating
recursively, proves that the integral of this word is \eqref{eq:MPL}.
For example,
\begin{equation}
 I(Z^3e_1e_1)=\Li_{4,1}(i,1),\qquad
 I(Z^3e_1e_3)=\Li_{4,1}(i,-1).
 \label{eq:examplewords}
\end{equation}
Using $e_2$ rather than $e_3$ as the final letter in the second expression
would give a different value. This is one of the convention errors that
an exact certificate must prevent.

\subsection{Convergence estimates used later}
For an admissible colored index word $B$ of depth $d$, let $B_N$ denote
its sum with outer index at most $N$. Its tail obeys
\begin{equation}
 B_N=B_\infty+O\!\left(\frac{(1+\log N)^d}{N}\right).
 \label{eq:tailB}
\end{equation}
For outer index at least two, comparison with
$\sum_{n>N}(1+\log n)^{d-1}/n^2$ proves the bound. For outer index one,
admissibility forces its first series color $c$ to be a nontrivial root
of unity. If $A_{n-1}$ is the inner partial sum, the coefficient
$b_n=A_{n-1}/n$ satisfies
\[
 b_n=O((1+\log n)^{d-1}/n),\qquad
 b_{n+1}-b_n=O((1+\log n)^{d-1}/n^2).
\]
The second estimate follows by separating the change in $1/n$ from the
new inner summand; the latter is bounded by a power of $\log n$ divided
by at least $n$. The partial sums of $c^n$ are bounded. Summation by parts
therefore gives \eqref{eq:tailB}. This also verifies ordinary convergence
and agreement with radial Abel limits.

\subsection{Shuffle, stuffle, and conjugation}
Let $u\sh v$ denote the shuffle of words, with multiplicity for repeated
letters. The recursion is
\begin{equation}
 (au)\sh(bv)=a\bigl(u\sh(bv)\bigr)+b\bigl((au)\sh v\bigr),
 \qquad u\sh\varnothing=u.
 \label{eq:shuffle}
\end{equation}
Partitioning the product of two ordered simplices gives
$I(u\sh v)=I(u)I(v)$ whenever both words are admissible.

The stuffle is formed in relative-color index coordinates. If
$A=(a,c)$ and $B=(b,d)$, put $A\oplus B=(a+b,c+d\bmod4)$ and define
\begin{equation}
 (AU)*(BV)=A\bigl(U*(BV)\bigr)+B\bigl((AU)*V\bigr)
             +(A\oplus B)(U*V).
 \label{eq:stuffle}
\end{equation}
Sorting finite summation indices gives this identity exactly before
passing to a limit. Define the raw word relation
\begin{equation}
 \mathcal R(u,v)=u\sh v-
 \word\bigl(\ind(u)*\ind(v)\bigr),
 \label{eq:rawDS}
\end{equation}
where the word map is extended linearly. Its integral is zero for
admissible $u,v$. Section~\ref{sec:regularization} proves the only
non-admissible-factor case used in the certificate.

Complex conjugation fixes $Z,e_0,e_2$ and exchanges $e_1,e_3$. For a word
polynomial $P$, define $\Delta P=P-\overline P$. The computational
encoding uses integers $-1,0,1,2,3$, where $-1$ means $Z$ and $j\ge0$
means $e_j$; in particular, the integer $0$ is \emph{not} the zero form.
Order encoded words lexicographically. The compressed odd projection is
\begin{equation}
 \Odd\Bigl(\sum_w c_ww\Bigr)
 =\sum_{w<\overline w}(c_w-c_{\overline w})[w].
 \label{eq:odd}
\end{equation}
Self-conjugate words disappear. Thus
\begin{equation}
 \Imag I(P)=\sum_{w<\overline w}
              (c_w-c_{\overline w})\Imag I(w).
 \label{eq:oddmeaning}
\end{equation}
There is no hidden factor of two in \eqref{eq:odd}; the factor appears
only when $P$ itself is already antisymmetric.

\section{The one-divergence double-shuffle lemma}
\label{sec:regularization}

\begin{lemma}[One leading $\Li_1(1)$ is sufficient]\label{lem:oneDiv}
Let $v$ be a nonempty admissible word. The formal polynomial
$\mathcal R(e_0,v)$ has only admissible terms after like words are
collected, and
\begin{equation}
 I\bigl(\mathcal R(e_0,v)\bigr)=0.
 \label{eq:oneDiv}
\end{equation}
No assertion about products with two divergent factors is required.
\end{lemma}
\begin{proof}
The only non-admissible word in either product expansion is $e_0v$.
It occurs with coefficient one on each side and cancels. Indeed, a
shuffle can begin with $e_0$ only by taking it from the first factor,
since the first letter of $v$ is not $e_0$; the remaining word must be
$v$. In the stuffle, precisely the term obtained by putting the
index-color pair $(1,0)$ before the index word of $v$ is divergent.
Merged first indices are at least two and hence admissible.

It remains to justify that the finite parts of the two divergent terms
agree. Put $B_N=\Li_{\ind(v)}^{(N)}$ and $B=B_\infty$. The finite cutoff
of $e_0v$ is
\begin{align}
 D_N&=\sum_{n=1}^N\frac{B_{n-1}}n
 =BH_N-\sum_{n=1}^N\frac{B-B_{n-1}}n,\notag\\
 D_N-BH_N&\longrightarrow
 C_B:=-\sum_{n=1}^{\infty}\frac{B-B_{n-1}}n.
 \label{eq:cutoffconstant}
\end{align}
The last series converges absolutely by \eqref{eq:tailB}. For $0<r<1$,
the radial version is
\begin{equation}
 I(e_0v;r)=-B\log(1-r)
 -\sum_{n=1}^{\infty}\frac{r^n(B-B_{n-1})}{n},
 \label{eq:Abelconstant}
\end{equation}
so its constant term after subtracting $-B\log(1-r)$ is the same $C_B$.

The finite stuffle identity uses the product $H_NB_N$.
Its difference from $H_NB$ tends to zero by \eqref{eq:tailB}.
For the shuffle identity at endpoint $r$, the corresponding product is
$-\log(1-r)I(v;r)$. The tail bound and summation by parts imply
\[
 I(v;r)-B=O\bigl((1-r)(1+|\log(1-r)|)^{d+1}\bigr),
\]
so replacing $I(v;r)$ by $B$ also changes the product by a quantity
tending to zero. Every remaining term in the two product expansions is
admissible and has its ordinary limit. Subtracting the common divergent
term and then its common constant proves \eqref{eq:oneDiv}.
\end{proof}

\begin{remark}
This elementary lemma is consistent with the regularization comparison
in the general double-shuffle theory \cite{Zhao2010}, but avoids appealing
to that machinery for an unexamined divergent case. Precisely 123 of the
852 double-shuffle rows used below have first factor $e_0$. The other
729 have two admissible factors. No row with two divergent factors is
admitted by either verifier.
\end{remark}

\section{Cayley duality with explicit admissibility}
\label{sec:cayley}

\begin{theorem}[The level-four Cayley involution]\label{thm:cayley}
Let $\phi(t)=(1-t)/(1+t)$ and let $\tau=-\phi^*$ act on the forms
\eqref{eq:forms}. Then
\begin{equation}
\begin{array}{c|ccccc}
 a&Z&e_0&e_1&e_2&e_3\\ \hline
 \tau(a)&e_0-e_2&Z+e_2&e_2-e_3&e_2&e_2-e_1
\end{array}
\label{eq:tau}
\end{equation}
For a word $w=a_1\cdots a_n$ define
\begin{equation}
 \cay(w)=\tau(a_n)\cdots\tau(a_1),
 \label{eq:D}
\end{equation}
with multilinear expansion and ordinary concatenation on the right.
Every term of $\cay(w)$ is admissible when $w$ is admissible, and
\begin{equation}
 I(w)=I(\cay(w)),\qquad \cay^2(w)=w.
 \label{eq:duality}
\end{equation}
The operator commutes with complex conjugation and preserves weight,
but not necessarily depth.
\end{theorem}
\begin{proof}
The identity $\phi'(t)=-2/(1+t)^2$ gives the five pullbacks directly.
For example,
\[
 -\phi^*Z=\frac{2\dd t}{1-t^2}=e_0-e_2,
 \qquad -\phi^*e_0=\frac{\dd t}{t(1+t)}=Z+e_2.
\]
For $e_1$, simplifying
$2i\dd t/((1+t)((1-i)+(1+i)t))$ gives $e_2-e_3$;
conjugation gives the $e_3$ row. The $e_2$ row is immediate.

In the ordered-simplex integral for $I(w)$, substitute
$t_j=\phi(u_{n+1-j})$. This reverses the order of the variables and
reverses their orientations. The resulting form at each position is
$-\phi^*\omega_{a_j}$, which is exactly \eqref{eq:D}. Convergence permits
the substitution by first restricting away from the endpoints and
then taking limits.

To check that no divergent transformed word has been introduced,
note that the old final letter lies in $\{e_0,e_1,e_2,e_3\}$.
Every term of its transform avoids $e_0$ in the new first position.
The old first letter lies in $\{Z,e_1,e_2,e_3\}$, and every term of its
transform is a nonzero letter in the new last position. Hence all
transformed words are admissible. Finally $\phi^2$ is the identity,
so $\tau^2$ and $\cay^2$ are identities. The coefficients and the map
$\phi$ are real, which proves compatibility with conjugation.
\end{proof}

\subsection{The particular relation needed at weight five}
Put
\begin{equation}
 A=Z^3e_1e_1.
 \label{eq:A}
\end{equation}
The entire additional symmetry input in the certificate is
\begin{equation}
 I(A)=I\bigl((e_2-e_3)^2(e_0-e_2)^3\bigr).
 \label{eq:seed}
\end{equation}
The powers in \eqref{eq:seed} are concatenation powers, not shuffle
powers. The right side expands into 32 words before possible collection,
all of depth five and all admissible. This is a direct specialization
of Theorem~\ref{thm:cayley}, not an experimentally inferred relation.

The transformation supplies an explanation for why a calculation organized
only around a small double-polylogarithm basket can miss a relation:
the relation can become apparent only after passing through higher-depth
words. This observation does not assert that any specific standard
relation ideal is incomplete. It describes the mechanism of the
certificate actually delivered here.

\section{The exact finite certificate and the proof of \texorpdfstring{$S_4$}{S4}}
\label{sec:certificate}

\subsection{The bridge from \texorpdfstring{$S_p$}{Sp} to colored doubles}
\begin{lemma}\label{lem:bridge}
For $p>1$,
\begin{equation}
 S_p=\Imag\bigl(\Li_{p,1}(i,1)+\Li_{p,1}(i,-1)\bigr).
 \label{eq:bridge}
\end{equation}
The same identity holds for $p=1$ by convergent Abel limits.
\end{lemma}
\begin{proof}
The inner coefficient of the sum of the two doubles is
$\sum_{k<m}(1+(-1)^k)/k$. If $m=2n+1$, this equals $H_n$,
because only $k=2j$ survive. Even $m$ have zero imaginary coefficient
from $i^m$, whereas $\Imag i^{2n+1}=(-1)^n$. Absolute convergence
justifies the selection for $p>1$. The index-one case follows by first
inserting a radial factor and then using convergence of the boundary
values.
\end{proof}

\subsection{A target polynomial with no transcendental coefficients}
Set
\begin{equation}
 B=Z^2e_1Ze_1,
 \qquad C=Z^3e_1e_3,
 \qquad V=Z^3e_1,
 \qquad q=e_1-e_3.
 \label{eq:targetletters}
\end{equation}
Here $I(B)=\Li_{3,2}(i,1)$ and $I(V)=\Li_4(i)$.
In the shuffle algebra define
\begin{equation}
 T=56\Delta C-408\Delta A-192\Delta B
       -19q^{\sh5}-112(\Delta V\sh e_2).
 \label{eq:T}
\end{equation}
Unlike the powers in \eqref{eq:seed}, $q^{\sh5}$ is the fivefold shuffle
power. All coefficients in $T$ are integers.

\begin{lemma}[Translation of the scalar identity]\label{lem:T}
Let
\begin{equation}
 R_4=7S_4-58g_{4,1}-24g_{3,2}
      -\frac{19\pi^5}{512}+14\beta(4)L.
 \label{eq:R4}
\end{equation}
Then $T$ is antisymmetric under conjugation and
\begin{equation}
 I(T)=16iR_4.
 \label{eq:IT}
\end{equation}
\end{lemma}
\begin{proof}
The single-value evaluations needed here are
\[
 I(q)=\Li_1(i)-\Li_1(-i)=\frac{\pi i}{2},\qquad
 I(e_2)=-L,\qquad I(\Delta V)=2i\beta(4).
\]
They follow from $\Li_1(z)=-\log(1-z)$ on the indicated segments and
the series for $\Li_4(i)$. Consequently the first three terms in
\eqref{eq:T} contribute
\[
 112i\Imag I(C)-816ig_{4,1}-384ig_{3,2},
\]
the fourth contributes $-19i\pi^5/32$, and the fifth contributes
$224i\beta(4)L$. Substitute \eqref{eq:bridge} with $p=4$ and collect
terms to obtain \eqref{eq:IT}. Each term of $T$ changes sign under
conjugation; in the fourth term the shuffle power is odd.
\end{proof}

\subsection{The certificate identity}
\begin{lemma}[Exact rational word certificate]\label{lem:certificate}
There are 852 pairs $(u_j,v_j)$ and rational coefficients $c_j$, recorded
in \path{data/S4_certificate.json}, such that
\begin{equation}
 \boxed{\displaystyle
 \Odd(T)=224\Odd\bigl(A-\cay(A)\bigr)
              +\sum_{j=1}^{852}c_j\Odd\bigl(\mathcal R(u_j,v_j)\bigr).}
 \label{eq:certificate}
\end{equation}
For each pair, $v_j$ is admissible and $u_j$ is either admissible or
$e_0$. All total weights are five. Exactly 541 pairs have word lengths
$(1,4)$ and 311 have lengths $(2,3)$; exactly 123 have $u_j=e_0$.
\end{lemma}

\begin{proof}[Finite exact verification]
The certificate is a list of rational numbers and finite words, not a
list of approximate special values. Each double-shuffle row is regenerated
by \eqref{eq:shuffle}--\eqref{eq:rawDS}; the Cayley row is regenerated
from \eqref{eq:tau}. Apply \eqref{eq:odd} and subtract the right side of
\eqref{eq:certificate} from its left side using rational arithmetic.
Every resulting coefficient is zero. The delivered primary verifier
\path{code/verify_s4.py} performs this calculation and rejects any
unallowed factor or non-admissible resulting word. Its receipt reports
zero nonzero residual coordinates.

A second verifier, \path{code/verify_s4_independent.py}, does not import
the primary algebra implementation. It enumerates shuffle placements
and ordered index merges, retains full words rather than compressed
odd coordinates, and checks
\begin{equation}
 2T=224\Delta\bigl(A-\cay(A)\bigr)
           +\sum_{j=1}^{852}c_j\Delta\mathcal R(u_j,v_j).
 \label{eq:fullcertificate}
\end{equation}
All full-word coefficients again vanish. The two checks establish
\eqref{eq:certificate} as a finite identity over $\Q$. The complete
coefficient data, the reconstruction rules, and executable verification
kernels are part of this proof package. Their use is essential; a claimed
rank or a numerical residual is not substituted for the coefficient
identity.
\end{proof}

\begin{proof}[Proof of Theorem~\ref{thm:S4}]
Each term $\mathcal R(u_j,v_j)$ has zero integral by ordinary
double-shuffle or Lemma~\ref{lem:oneDiv}. The Cayley row has zero integral
by Theorem~\ref{thm:cayley}. Taking imaginary parts in
\eqref{eq:certificate}, using \eqref{eq:oddmeaning}, therefore yields
$\Imag I(T)=0$. Lemma~\ref{lem:T} gives $16R_4=0$. Dividing
\eqref{eq:R4} by seven proves \eqref{eq:S4}, and subtracting $g_{4,1}$
from \eqref{eq:bridge} proves \eqref{eq:mixed41}.
\end{proof}

\subsection{Size, provenance, and limitations of the finite computation}
The sparse certificate has 853 nonzero coefficients including its one
Cayley row. The largest numerator and denominator have 32 and 24 decimal
digits, respectively. The target has 24 compressed odd coordinates.
For orientation, there are 2000 admissible length-five words over this
alphabet. Of them, 108 are fixed by conjugation, leaving 946 possible
compressed odd coordinates. These counts concern words, not independent
periods.

The candidate search generated weight splits $(1,4)$ and $(2,3)$ at all
admissible depths, allowing the one exceptional first factor $e_0$.
After removal of duplicate primitive odd rows there were 1380
nonzero double-shuffle rows. A modular calculation was useful for
locating a candidate; exact rational elimination then produced the
certificate. Neither modular rank evidence nor the elimination software
is needed to replay the proof. The verifiers consume only the fixed
certificate and the elementary algebraic rules. The optional script
\path{code/regenerate_s4.py} reconstructs the prescribed candidate rows
and performs the exact elimination anew; its output is then checked by
both ordinary verifiers.

No minimality statement is attached to the 852 rows, their coefficient
heights, or the additional Cayley relation. In particular, this is not a
proof that a suitably enlarged standard-relation system could never
recover the same identity without an explicit symmetry row. A much
shorter hand proof would be a valuable improvement.

\section{Further proved special-value identities}
\label{sec:consequences}

\begin{corollary}[A mixed-color antisymmetric reduction]\label{cor:antisym}
One has
\begin{align}
 &\Imag\bigl(\Li_{4,1}(i,-1)-\Li_{1,4}(-1,i)\bigr)\notag\\
 &\hspace{8mm}=\frac{102g_{4,1}+48g_{3,2}}7
       +\frac{79\pi^5}{10752}-3\beta(4)L.
 \label{eq:antisymmetric}
\end{align}
\end{corollary}
\begin{proof}
The ordinary stuffle product gives
\[
 \Li_{4,1}(i,-1)+\Li_{1,4}(-1,i)
  =\Li_4(i)\Li_1(-1)-\Li_5(-i).
\]
Its imaginary part is $-\beta(4)L+\beta(5)$. The classical odd-beta
evaluation gives $\beta(5)=5\pi^5/1536$. Subtract this sum from twice
\eqref{eq:mixed41}. The rational coefficient is
$19/1792-5/1536=79/10752$, proving \eqref{eq:antisymmetric}.
\end{proof}

\begin{corollary}[An integral evaluation]\label{cor:integral}
The real logarithmic integral
\begin{equation}
 J_4=\int_0^1\frac{(-\log x)^3\log(1+x^2)}{1+x^2}\dd x
 \label{eq:J4}
\end{equation}
satisfies
\begin{equation}
 J_4=-\frac{348g_{4,1}+144g_{3,2}}7
       -\frac{57\pi^5}{1792}+12\beta(4)L.
 \label{eq:J4value}
\end{equation}
\end{corollary}
\begin{proof}
The generating series
$\sum_{n\ge0}H_nz^n=-\log(1-z)/(1-z)$ implies
$\log(1+x^2)/(1+x^2)=-\sum_{n\ge0}(-1)^nH_nx^{2n}$.
After multiplying by $(-\log x)^3$ and integrating, absolute
summability of the integrated terms gives $J_4=-6S_4$.
Now apply \eqref{eq:S4}.
\end{proof}

The same argument gives the integral representation used in the numerical
checks:
\begin{equation}
 S_p=-\frac1{(p-1)!}\int_0^1
       \frac{(-\log x)^{p-1}\log(1+x^2)}{1+x^2}\dd x.
 \label{eq:Spintegral}
\end{equation}
For $a,b\ge1$, another useful representation is
\begin{equation}
 g_{a,b}=\Imag\frac1{(a-1)!}\int_0^1
       \frac{(-\log x)^{a-1}i\Li_b(ix)}{1-ix}\dd x.
 \label{eq:gintegral}
\end{equation}
It follows by summing the inner geometric series before applying the
Mellin integral for $n^{-a}$. For $a\ge2$ and $b=1$, the identity
$\Li_{1,1}(z,1)=\tfrac12\log^2(1-z)$ gives the independent real form
\begin{equation}
 g_{a,1}=-\frac1{2(a-2)!}\int_0^1
  \frac{(-\log x)^{a-2}}x\log(1+x^2)\arctan x\dd x.
 \label{eq:g1integral}
\end{equation}
These integrals avoid a slowly convergent direct nested summation.
Their floating-point quadrature is used for discovery only; the
certification below uses exact rational sums instead.

\section{Signed moments for affine harmonic coefficients}
\label{sec:moments}

\subsection{A general Gamma-density bound}
The ordinary alternating-series remainder theorem cannot simply be
applied after arbitrary transformations of harmonic coefficients.
A signed-moment representation gives a different, uniform control.
For a finite signed measure $\nu$, our convention is
$\norm\nu_{\TV}=|\nu|([0,1])$, not half this quantity.

\begin{theorem}[Affine harmonic signed moments]\label{thm:affine}
Let $a,r,d$ be positive integers and let $b>0$ be real. Set
\begin{equation}
 a_k=\frac{H_{rk}^{(b)}}{(dk+1)^a},\qquad k\ge0,
 \label{eq:affineseq}
\end{equation}
and define
\begin{equation}
 m_1=1,\qquad
 m_a=\frac{(a-1)^{a-1}e^{-(a-1)}}{\Gamma(a)}\quad(a>1).
 \label{eq:ma}
\end{equation}
There is a finite signed measure $\nu$ on $[0,1]$, with no atom at one,
such that $a_k=\int_0^1x^k\dd\nu(x)$ and
\begin{equation}
 \norm\nu_{\TV}\le\frac{2r}{d}\,m_a\,b\zeta(b+1).
 \label{eq:TVsharp}
\end{equation}
For $b>1$ one also has
\begin{equation}
 \norm\nu_{\TV}\le2\zeta(b).
 \label{eq:TVb}
\end{equation}
In particular, for integer $b\ge1$, usable rational bounds are
\begin{equation}
 \norm\nu_{\TV}\le
 \begin{cases}
 10r/(3d),&b=1,\\
 10/3,&b\ge2.
 \end{cases}
 \label{eq:TVcoarse}
\end{equation}
\end{theorem}
\begin{proof}
Let $Y$ have the Gamma density $y^{a-1}e^{-y}/\Gamma(a)$ and let
$\mu$ be the law of $X=e^{-dY}$. Its moments are
\begin{equation}
 \int x^k\dd\mu(x)=\mathbb E(e^{-dkY})=(dk+1)^{-a}.
 \label{eq:mu}
\end{equation}
In the logarithmic coordinate $s=-\log x$, the density is
\[
 f(s)=\frac{s^{a-1}e^{-s/d}}{d^a\Gamma(a)}\,\mathbf 1_{s>0}.
\]
Its maximum is $m_a/d$. Extended by zero to the negative half-line,
its distributional derivative has total variation $2m_a/d$.
For $a>1$, this is the sum of the rise to the maximum and the fall
back to zero. For $a=1$, the jump at zero contributes $1/d$ and
the decreasing exponential contributes another $1/d$.

The $L^1$ distance between an integrable function of bounded variation
and its translate by $h\ge0$ is at most $h$ times the variation of its
distributional derivative. This follows, for a smooth function, by
integrating $f(s-h)-f(s)=-\int_{s-h}^s f'(v)\dd v$ and applying
Fubini; approximation by smooth convolutions proves the bounded-variation
case. Consequently, if $D_\rho\mu$ denotes pushforward under
$x\mapsto\rho x$, then
\begin{equation}
 \norm{\mu-D_{e^{-rt}}\mu}_{\TV}
 \le\min\left(2,\frac{2m_ar}{d}t\right),\qquad t>0.
 \label{eq:translation}
\end{equation}

The elementary integral for the generalized harmonic number is
\begin{equation}
 H_{rk}^{(b)}=\frac1{\Gamma(b)}\int_0^\infty
       \frac{t^{b-1}(1-e^{-rkt})}{e^t-1}\dd t.
 \label{eq:Hlaplace}
\end{equation}
To prove it, expand the finite geometric expression
$(1-e^{-rkt})/(e^t-1)=\sum_{j=1}^{rk}e^{-jt}$ and integrate termwise.
Define the signed measure by a total-variation-convergent integral,
\begin{equation}
 \nu=\frac1{\Gamma(b)}\int_0^\infty
       \frac{t^{b-1}}{e^t-1}
       \bigl(\mu-D_{e^{-rt}}\mu\bigr)\dd t.
 \label{eq:nu}
\end{equation}
Near zero, \eqref{eq:translation} supplies an additional factor $t$,
so the norm of the integrand is $O(t^{b-1})$, integrable for every
$b>0$. At infinity the exponential denominator controls it. This
justifies the measure integral, Fubini for bounded moment functions,
and the absence of an atom at one. Combining \eqref{eq:mu} and
\eqref{eq:Hlaplace} gives exactly the moments \eqref{eq:affineseq}.

Using the linear bound in \eqref{eq:translation} yields
\[
 \norm\nu_{\TV}\le
 \frac{2m_ar}{d\Gamma(b)}\int_0^\infty\frac{t^b}{e^t-1}\dd t
 =\frac{2r}{d}m_ab\zeta(b+1),
\]
which proves \eqref{eq:TVsharp}. For $b>1$, use instead the constant
bound two in \eqref{eq:translation}, giving \eqref{eq:TVb}.
For integer $a$, $m_a\le1$: with $n=a-1\ge1$, this is the elementary
inequality $e^n\ge n^n/n!$. Finally $\zeta(2)<5/3$, for instance by
summing $n^{-2}$ through $n=5$ and bounding the remaining tail by
$\int_5^\infty x^{-2}\dd x$. Since $\zeta(b)\le\zeta(2)$ for
integer $b\ge2$, \eqref{eq:TVcoarse} follows.
\end{proof}

\begin{remark}[Refinements and comparison with the source kernel]
The full bound obtained in the proof is
\begin{equation}
 \norm\nu_{\TV}\le\frac2{\Gamma(b)}\int_0^\infty
 \frac{t^{b-1}}{e^t-1}\min\left(1,\frac{m_ar}{d}t\right)\dd t.
 \label{eq:TVmin}
\end{equation}
It can improve either elementary estimate. The factor $m_a$ in
\eqref{eq:TVsharp} decreases on the scale $a^{-1/2}$, by Stirling's
formula. The proof therefore provides an outer-index-sensitive bound,
not just a constant independent of the indices.

When $r=d$, the signed kernel already constructed in the inspected
manuscript for $H_{n-1}^{(b)}/n^a$ gives another route, by pushforward
under $u\mapsto u^d$. Its base positive and negative terms each have
integral one and its recursion is an $L^1$ contraction, so it gives
the additional coarse bound two. We do not treat this special-case
kernel as new. Theorem~\ref{thm:affine} separates $r$ and $d$ and
supplies the explicit Gamma-density refinement. The delivered interval
receipts deliberately use the simpler bounds \eqref{eq:TVcoarse}, so
neither a numerical value of $e$ nor a Gamma evaluation is trusted by
the proof code.
\end{remark}

\subsection{Exact Euler acceleration for signed measures}
\begin{theorem}[Certified Euler transform]\label{thm:Euler}
Let $a_k=\int_0^1x^k\dd\nu(x)$, where $\nu$ is a finite signed measure
with $\nu(\{1\})=0$ and $\norm\nu_{\TV}\le C$. Then the ordinary
alternating series converges, and for every $N\ge1$,
\begin{equation}
 \sum_{k=0}^{\infty}(-1)^ka_k=E_N+\eps_N,
 \qquad |\eps_N|\le C2^{-N},
 \label{eq:Eulerbound}
\end{equation}
where, with the decrement difference $\delta a_k=a_k-a_{k+1}$,
\begin{align}
 E_N&=\sum_{j=0}^{N-1}\frac{\delta^ja_0}{2^{j+1}}\notag\\
 &=2^{-N}\sum_{k=0}^{N-1}(-1)^k
       \left(2^N-\sum_{j=0}^k\binom Nj\right)a_k.
 \label{eq:Eulertailweights}
\end{align}
The sign convention for $\delta$ is opposite to the forward difference
used in some statements of Euler's transformation.
\end{theorem}
\begin{proof}
The finite alternating kernel $\sum_{k=0}^M(-x)^k$ is bounded in
absolute value by two on $[0,1]$, and converges to $(1+x)^{-1}$ for
$x<1$. Dominated convergence with respect to $|\nu|$, using the
absence of an atom at one, proves ordinary convergence and the formula
\[
 \sum_{k\ge0}(-1)^ka_k=\int_0^1\frac{\dd\nu(x)}{1+x}.
\]
The finite geometric identity
\begin{equation}
 \frac1{1+x}=\sum_{j=0}^{N-1}\frac{(1-x)^j}{2^{j+1}}
      +\frac{(1-x)^N}{2^N(1+x)}
 \label{eq:kernelEuler}
\end{equation}
gives \eqref{eq:Eulerbound}, because the last kernel is bounded by
$2^{-N}$. Expanding $(1-x)^j$ gives the first formula for $E_N$.
Its coefficient of $a_k$ is
$(-1)^k\sum_{j=k}^{N-1}2^{-j-1}\binom jk$.
The binomial-tail identity
\[
 \sum_{j=k}^{N-1}2^{-j-1}\binom jk
   =2^{-N}\sum_{j=k+1}^{N}\binom Nj
\]
follows by induction on $N$ using Pascal's identity, and proves the
second formula in \eqref{eq:Eulertailweights}.
\end{proof}

No complete monotonicity of the harmonic sequence has been assumed.
Its representing measure is signed and has total mass zero because
$a_0=0$. The total-variation bound, rather than positivity of finite
differences, is what certifies the transform.

\subsection{A rational-only implementation for all constants used here}
The imaginary-part selection in \eqref{eq:MPL} gives
\begin{equation}
 g_{a,b}=\sum_{k=0}^\infty
       \frac{(-1)^kH_{2k}^{(b)}}{(2k+1)^a}.
 \label{eq:gseries}
\end{equation}
Apply Theorem~\ref{thm:affine} with $(r,d)=(1,2)$ for $S_p$ and
$(r,d)=(2,2)$ for $g_{a,b}$. The rational constants actually used are
$C=5/3$ for $S_p$ and $C=10/3$ for each $g_{a,b}$.

The sequences $(2k+1)^{-s}$ and $(k+1)^{-s}$ are positive moment
sequences of mass one by the same Gamma integral. Thus the beta and
eta values have Euler error at most $2^{-N}$. Finally,
\begin{equation}
 \pi=4\beta(1),\qquad L=\eta(1),\qquad
 \zeta(s)=\frac{\eta(s)}{1-2^{1-s}}\quad(s>1).
 \label{eq:singleEuler}
\end{equation}
Every center in \eqref{eq:Eulertailweights} and every chosen error
radius is rational for the integer indices used in this paper.

The implementation maintains the binomial tail
$2^N-\sum_{j=0}^k\binom Nj$ incrementally. It needs $O(N)$ rational
arithmetic operations for the transform once the coefficient list is
available; producing the harmonic list needs $O(rN)$ rational additions.
These are arithmetic-operation counts, not linear bit-complexity claims.
The numerator and denominator sizes grow with $N$.

Exact interval arithmetic encloses products and powers without treating
reused constants as independent random errors. The code multiplies
interval endpoints and takes their minimum and maximum, so overestimation
from dependency is harmless. To keep stored endpoints compact, it rounds
outward to rationals with denominator $10^{350}$, again using only integer
arithmetic. Decimal strings in the receipts are display summaries;
the rational lower and upper endpoints are the authoritative data.

\section{A weight-seven conjecture and a rejected weight-nine relation}
\label{sec:discovery}

\subsection{The surviving \texorpdfstring{$S_6$}{S6} candidate}
\begin{conjecture}[A mixed Gaussian identity of weight seven]\label{conj:S6}
With the notation \eqref{eq:Sp}--\eqref{eq:notation},
\begin{equation}
 \boxed{\begin{aligned}
 S_6={}&-\frac{12334}{527}g_{6,1}-\frac{384}{31}g_{5,2}
        -\frac{2136}{527}g_{4,3}\\
        &+\frac{3179\pi^7}{5713920}
        -2\beta(6)L-\frac{801}{2108}\beta(4)\zeta(3).\end{aligned}}
 \label{eq:S6}
\end{equation}
\end{conjecture}

The discovery search used the ordered basket
\begin{equation}
 (S_6,g_{6,1},g_{5,2},g_{4,3},\pi^7,
          \beta(6)L,\beta(2)\zeta(5),\beta(4)\zeta(3)).
 \label{eq:S6basket}
\end{equation}
At 110 working decimal digits, PSLQ with tolerance $10^{-95}$,
coefficient ceiling $10^{12}$, and at most 10000 iterations returned
\begin{equation}
\begin{split}
 (&-97136640,-2273402880,-1203240960,-393707520,\\
  &\hspace{9mm}54043,-194273280,0,-36910080).
\end{split}
\label{eq:S6vector}
\end{equation}
This primitive integer vector is equivalent to \eqref{eq:S6}; its
zero coefficient on $\beta(2)\zeta(5)$ is part of the observation,
not a pre-imposed omission. The decimal data are retained in the
package. A separate quadrature at 180 digits supported the relation.
The exact rational check below is more informative than agreement of
two floating-point evaluations.

Let $R_6$ be the left side minus the right side of \eqref{eq:S6}.
Using $N=1000$ in the certified Euler transform yields, after outward
rounding,
\begin{equation}
 -1.357\cdot10^{-299}<R_6<1.515\cdot10^{-299}.
 \label{eq:S6enclosure}
\end{equation}
The full rational interval has width less than
$2.871\cdot10^{-299}$. This is a proved enclosure for the residual,
not a proof that the residual is zero. In particular, it does not license
adding \eqref{eq:S6} as a relation to an exact word certificate.

For orientation, the underlying values start as follows; the displayed
decimals are rounded, not interval endpoints:
\begin{center}
\begin{tabular}{lr}
\toprule
Quantity&Approximate value\\
\midrule
$S_6$&$-0.0012883399816434530677630587571690$\\
$g_{6,1}$&$-0.0019411939935331449052635479564949$\\
$g_{5,2}$&$-0.0047581161681305610395611455068384$\\
$g_{4,3}$&$-0.0123729114645383875464708307852280$\\
\bottomrule
\end{tabular}
\end{center}
All monomials in \eqref{eq:S6} have total weight seven. The form is
therefore compatible with the weight grading. Compatibility and
high-precision agreement are useful selection criteria, but neither
substitutes for the missing exact reduction.

\subsection{A false candidate that looked convincing at discovery precision}
A broader search at weight nine used
\begin{equation}
\begin{split}
 \mathbf b_8=(&S_8,g_{8,1},g_{7,2},g_{6,3},g_{5,4},\pi^9,\\
            &\beta(8)L,\beta(2)\zeta(7),\beta(4)\zeta(5),
                         \beta(6)\zeta(3)).
\end{split}
\label{eq:S8basket}
\end{equation}
At 110 digits, with tolerance $10^{-90}$ and coefficient ceiling
$10^{12}$, a search returned
\begin{equation}
\begin{split}
 \mathbf c_8=(&-6451250972,-4072849461,-710107844,\\
             &-10984374763,-6393592639,31875,\\
             &-1955138863,3459066102,-1034180101,-1477786143).
\end{split}
\label{eq:S8vector}
\end{equation}
This is \emph{not} proposed as a conjecture. Define
$R_8=\mathbf c_8\cdot\mathbf b_8$.

\begin{proposition}[Certified rejection of the weight-nine near-relation]
\label{prop:S8false}
The explicit vector \eqref{eq:S8vector} is not an exact relation:
\begin{equation}
 -1.3274179020580459\cdot10^{-86}
 <R_8<
 -1.3274179020580457\cdot10^{-86}.
 \label{eq:S8exclusion}
\end{equation}
In particular, $R_8<0$.
\end{proposition}
\begin{proof}
Apply Theorems~\ref{thm:affine} and \ref{thm:Euler} with $N=1000$ to
each harmonic sum and \eqref{eq:singleEuler} to each single value.
Form the indicated products and linear combination in exact interval
arithmetic. The resulting rational interval, saved in
\path{data/S8_exclusion_N1000.json}, has width below
$2.080\cdot10^{-290}$ and is wholly contained in the interval
\eqref{eq:S8exclusion}. Both comparisons are exact comparisons of
rational endpoints. Hence zero is excluded.
\end{proof}

A separate search using 180-digit inputs, tolerance $10^{-160}$,
coefficient ceilings up to $10^{12}$ and at most 15000 iterations
returned no relation in this basket. That search failure does not
prove the absence of another relation, even below a height ceiling,
unless the algorithm supplies and one verifies the relevant exclusion
certificate. We assert only the observed failure and the rigorous
rejection of the particular vector \eqref{eq:S8vector}.

The numerical warning is not that PSLQ \cite{PSLQ} is unreliable as an
integer-relation algorithm. A tolerance applied to scaled coordinates
can produce a vector whose \emph{unnormalized} residual is not as small
as a casual reading of the tolerance suggests. The coefficient sizes,
normalizations, input errors, and recomputed residual all matter. The
rational interval test checks the actual proposed identity and makes
that distinction explicit.

\subsection{Summary of the three residual checks}
\begin{center}
\begin{tabular}{@{}p{.15\textwidth}p{.25\textwidth}p{.20\textwidth}p{.31\textwidth}@{}}
\toprule
Residual&Enclosure width, upper bound&Contains zero?&Conclusion\\
\midrule
$R_4$, \eqref{eq:R4}&$7.109\cdot10^{-299}$&Yes&Consistent with the separate exact word proof.\\[4pt]
$R_6$, \eqref{eq:S6}&$2.871\cdot10^{-299}$&Yes&Strong certified numerical evidence; still conjectural.\\[4pt]
$R_8$, \eqref{eq:S8vector}&$2.080\cdot10^{-290}$&No&The displayed candidate is rigorously false.\\
\bottomrule
\end{tabular}
\end{center}
The identity proof and the error certification are separate logical
objects. A vanishing word polynomial proves equality for exact periods;
a narrow interval containing zero does not. An interval excluding zero,
however, is enough to refute a proposed equality.

\section{Targeted manuscript corrections and integration}
\label{sec:corrections}

\subsection{Upgrade a conjecture, not an old proof}
The first integration change is a status upgrade. At the pinned snapshot,
\texttt{gaussian:conj:S4} is explicitly a conjecture, supported by numerical
integrals and not included in the exact shuffle certificate. That caution
was appropriate. Replace its conjecture environment by a theorem and
attach the Cayley proof or a precise reference to this complete proof
package. The old label may be retained for compatibility even though its
name contains \texttt{conj}.

The original-coordinate version in the source,
\begin{equation}
 S_4=\frac{4g_{4,1}-3g_{3,2}-9g_{2,3}}7
       +\frac{\pi^5}{224}-\frac{27\beta(2)\zeta(3)}{224}
       -2\beta(4)L,
 \label{eq:S4old}
\end{equation}
is the same identity. The source already proves
\[
 g_{2,3}=-6g_{4,1}-3g_{3,2}-\frac{\pi^5}{1536}
                         -\frac{3\beta(2)\zeta(3)}{32}.
\]
Substitution in \eqref{eq:S4old} gives \eqref{eq:S4} exactly. Thus both
versions can be promoted, without re-proving or altering the existing
two-shuffle elimination.

The nearby sentence saying that no functional-equation proof is supplied
must be replaced, and the open-problem summary in
\path{chapters/10-discovery.tex} must no longer list this particular
$S_4$ identity as unresolved. Historical drafts should remain untouched
as provenance; only the canonical status should change.

\subsection{Qualify a non-reducibility statement}
In the subsection ``Family 2: antisymmetric reductions to the $(i,1)$
algebra'' in \path{chapters/04-depth.tex}, the prose says that at higher
weight the antisymmetric parts do not reduce against the $(i,1)$ basket
alone. Read as a universal nonmembership statement, this is too strong
for the search evidence described there. It is now also directly
qualified by the explicit weight-five mixed antisymmetric reduction
\eqref{eq:antisymmetric}.

A suitable replacement is:
\begin{quote}
The stated searches did not recover all higher-weight antisymmetric
reductions in the restricted basket. This is a limitation of those
searches, not a proof of nonmembership. In particular, the Cayley
certificate gives the weight-five mixed reduction displayed below.
\end{quote}
This correction does not change the valid formal-quotient counts in
other parts of the source. Those counts must continue to be described
as dimensions of the specified relation systems, not as proved
arithmetic dimensions.

\subsection{What should remain explicitly unresolved}
Do not replace $S_4$ by an assertion that all even-index $S_{2m}$ have a
known reduction. The status of $S_6$ here is conjectural, and the naive
weight-nine search provided a false near-relation. The manuscript's
caveats about minimal depth and numerical independence remain necessary.
Similarly, no general completeness result for level-four relations is
established by this certificate.

The package supplies a replacement subsection, a bibliography entry,
and label-anchored editing notes under \path{integration/}. No changes
have been pushed to the repository. The article and certificates are
intended to be reviewed and integrated without overwriting historical
reports or their validation receipts.

\section{Further research questions}
\label{sec:future}

\subsection{Compress the exact \texorpdfstring{$S_4$}{S4} certificate}
The most immediate problem is to replace 852 double-shuffle rows by a
small human-readable chain. The single Cayley seed already identifies
a geometrically meaningful source of the relation. A useful intermediate
goal is a certificate grouped into a few products of lower-weight
identities, with all cancellations visible. Such a compression should
be measured by row count, coefficient height, and maximum intermediate
depth, not merely by a short final formula.

A systematic route is to reduce the 32 depth-five words in
\eqref{eq:seed} in a structured basis before eliminating depth-two
coordinates. An exact sparse optimization or a carefully chosen
Lyndon-basis normal form could produce a much smaller certificate.
Any candidate compressed proof can be checked against the already
verified target without recomputing period values.

\subsection{Prove the \texorpdfstring{$S_6$}{S6} identity}
Conjecture~\ref{conj:S6} is the highest-priority new special-value target.
Its coefficients are fixed, its residual is rigorously tiny, and the
$S_4$ proof supplies a plausible mechanism. One may test Cayley images
of $Z^5e_1e_1$ and other weight-seven Gaussian double words, together
with products and lower-weight symmetry relations. The objective is
an exact rational certificate, not another increment of floating-point
precision.

The unexpectedly absent $\beta(2)\zeta(5)$ term suggests an elimination
pattern worth understanding, but does not itself justify an all-weight
formula. A structured proof should explain the denominator factors
$17$ and $31$ that occur in the reduced coefficients, or show that they
are merely basis artifacts.

\subsection{Determine the correct weight-nine basket}
The rejection of \eqref{eq:S8vector} leaves several logically distinct
possibilities: there may be a genuine relation in the same basket with
larger coefficients; additional lower-weight products may be required;
or further depth-two or higher-depth coordinates may survive in the
chosen relation quotient. Finite exact elimination, with an explicitly
specified relation ideal, can distinguish some of these possibilities
at the formal level. It cannot by itself prove transcendence or
arithmetic independence of the evaluated periods.

A useful deliverable would be a complete weight-nine relation search
with a reproducible basis specification, exact rank certificates,
and a separate ledger of which reductions are numerical, formal,
or analytic. Missing a relation at one precision should never be
converted into a dimension assertion.

\subsection{Compare saturated relation systems}
The 1380-row search used here is deliberately limited: products of
weights $(1,4)$ and $(2,3)$, with only one-divergence regularization.
A fair comparison with ``standard relations'' in the broader literature
must include the chosen distribution relations, admissible lifts of
lower-weight identities, and weight-one seeds. Determine whether the
specific Cayley relation lies in the saturated ideal so obtained,
and record both the exact quotient and the explicit membership or
nonmembership certificate in that finite model.

At higher weights, distinguish adding new geometric seed relations from
multiplying already known seeds by lower-weight words. This may make
symmetry-based certificates far smaller and avoid rediscovering the
same structural relation in many disguises.

\subsection{Improve certified arithmetic with the refined moment bounds}
Theorem~\ref{thm:affine} gives the factor $m_a$ and the sharper integral
\eqref{eq:TVmin}, but the current replay uses coarse rational constants.
Implementing rational upper bounds for $m_a$ via finite lower bounds
for $e^{a-1}$ would preserve exact arithmetic while reducing the
required term count. The resulting evaluator should be compared with
the manuscript's Chebyshev and H\"older methods in bit operations,
memory, and certification overhead, not just in nominal geometric
rate.

Another natural extension is to products of harmonic numbers. Products
of signed moment measures correspond to multiplicative convolution;
this gives an immediate bound by the product of total variations,
but may be far from sharp because of cancellation. Whether the
Gamma-translation approach gives useful non-product bounds is a
concrete analytic question.

\subsection{Formalize the small trusted core}
A proof-assistant implementation could first formalize the finite word
identity, leaving the analytic interpretation as a separately proved
layer. The data are finite lists of integers and rationals, so this
stage is substantially smaller than formalizing a general multiple
polylogarithm theory. A complete formalization would additionally need
the iterated-integral shuffle, the convergent finite-stuffle limit,
the one-divergence lemma, and the Cayley substitution with endpoint
control.

The two existing Python verifiers are useful independent checks, but
they are not a substitute for such kernel-level formal verification.
Their agreement reduces implementation risk without changing that
logical distinction.

\section{Conclusion}
The inspected mixed Gaussian $S_4$ conjecture is resolved by an explicit
exact certificate whose only additional geometric input is one Cayley
identity. The proof changes the status of a concrete manuscript formula,
produces a mixed antisymmetric reduction, and exhibits how passing
through higher-depth words can expose a relation invisible to a
restricted search.

The continuation also supplies a carefully separated experimental
frontier. The $S_6$ identity survives a rational error certificate at
approximately 299 decimal places, whereas an apparently persuasive
$S_8$ candidate is rigorously rejected. These are different mathematical
outcomes and are recorded as such. The complete package preserves the
finite proof data, independent verification, error bounds, and
integration notes needed to extend the work without confusing numerical
evidence with exact identities.

\appendix
\section{Certificate format and the verification kernel}
\label{app:certificate}

The JSON schema is \texttt{proveit.cayley-s4.word-certificate.v1}. A
nonzero row has one of two forms:
\begin{lstlisting}
{"kind": "double_shuffle", "u": [ ... ], "v": [ ... ],
 "coefficient": "numerator/denominator"}
{"kind": "cayley", "word": [-1,-1,-1,1,1],
 "coefficient": "224"}
\end{lstlisting}
The coefficient multiplies the \emph{raw} row \eqref{eq:rawDS}, not
its primitive normalization. This convention is important because the
search initially normalized duplicate rows by signed integer gcds;
those normalizations have been absorbed into the published rational
coefficients. The TSV file is an equivalent human-readable listing.

The essential verification loop is small:
\begin{lstlisting}[language=Python]
residual = odd_projection(target())
for row in certificate["rows"]:
    c = Fraction(row["coefficient"])
    if row["kind"] == "double_shuffle":
        raw = double_shuffle(tuple(row["u"]), tuple(row["v"]))
    elif row["kind"] == "cayley":
        w = tuple(row["word"])
        raw = plus({w: 1}, cayley(w), -1)
    else:
        raise ValueError("Unknown relation type")
    residual = plus(residual, odd_projection(raw), -c)
if residual:
    raise ArithmeticError("Nonzero exact word residual")
\end{lstlisting}
The delivered implementation additionally validates total weight,
admissibility, duplicate identifiers, and nonzero coefficients. The
complete word-algebra recursions are in \path{code/word_algebra.py}.
The alternative verifier uses placement subsets for shuffles and an
explicit work list of merge paths for stuffles; it checks full
antisymmetrization rather than reusing \texttt{odd\_projection}.

The certificate SHA-256 digest is
\begin{center}\small
\texttt{4fd9d35fe6d1e343451034a6a35af890}\\
\texttt{07633e412a8bdafa91cb66f1dce769ed}.
\end{center}
Concatenate the two displayed lines without whitespace. The digest
identifies the exact JSON artifact verified here; it does not certify
mathematical correctness by itself.

\section{Reproduction and artifact map}
\label{app:repro}

The exact word and interval checks require only Python 3.10 or newer and
its standard library. The executed environment is recorded in
\path{logs/environment.json}. Optional quadrature and PSLQ replay uses
\texttt{mpmath}; neither it nor a linear solver is trusted by the exact
verifiers. From the package root, run:
\begin{lstlisting}
python code/verify_s4.py
python code/verify_s4_independent.py
python code/test_kernels.py
python code/verify_receipts.py
\end{lstlisting}
The last command recomputes all three $N=1000$ rational enclosures and
compares their endpoints exactly, without overwriting the recorded files.
It also verifies the numerical inequalities printed in this paper.
The optional \path{code/regenerate_s4.py} reconstructs the certificate
by exact elimination. Complete command lines for individual enclosures,
regeneration, and exploratory PSLQ searches are in \path{README.md}.

To rebuild the article with three serial PDFLaTeX passes, run
\begin{lstlisting}
python code/build_article.py
\end{lstlisting}
The supplied PDF was compiled and rendered for review. Rebuilt PDFs can
differ in metadata and need their own visual review. The manifest pins
the distributed artifacts; it is not a promise of byte-identical PDF
builds. In particular, a changed PDF is not covered by the original
visual-review receipt.

\begin{center}
\small
\begin{tabular}{@{}p{.43\textwidth}p{.51\textwidth}@{}}
\toprule
Artifact&Purpose\\
\midrule
\path{article/cayley_s4.tex}, PDF&Article source and reading copy.\\[3pt]
\path{data/S4_certificate.json}, TSV&Complete rational proof certificate.\\[3pt]
\path{code/verify_s4*.py}&Two exact word-certificate replays.\\[3pt]
\path{code/certified_euler.py}&Rational centers, analytic radii, interval arithmetic.\\[3pt]
\path{data/*N1000.json}&Exact rational residual enclosures.\\[3pt]
\path{data/*quadrature*}, discovery data&Non-certified experiment records.\\[3pt]
\path{integration/}&Replacement subsection and corrections.\\[3pt]
\path{logs/}, \path{MANIFEST.sha256}&Executed checks, environment, and hashes.\\
\bottomrule
\end{tabular}
\end{center}

\begingroup
\small
\begin{thebibliography}{9}
\setlength{\itemsep}{4pt}
\raggedright
\bibitem{ProveIt}
ProveIt contributors,
\emph{Polylogarithms and their Arithmetic Bridges}, collective manuscript,
\path{Analysis/Polylogarithms/docs/manuscript}, ProveIt repository.
Snapshot \texttt{a2a4cf58c49c745058c40e4a6748d472a3420f18}.
\href{https://github.com/VladimirReshetnikov/ProveIt/tree/a2a4cf58c49c745058c40e4a6748d472a3420f18/Analysis/Polylogarithms/docs/manuscript}{Pinned repository snapshot and canonical manuscript}.

\bibitem{Zhao2008}
J.~Zhao, Multiple polylogarithm values at roots of unity,
\emph{C.~R. Math. Acad. Sci. Paris} \textbf{346} (2008), 1029--1032.
\href{https://doi.org/10.1016/j.crma.2008.09.011}{doi:10.1016/j.crma.2008.09.011}.
\href{https://arxiv.org/abs/0810.1064}{arXiv:0810.1064}.

\bibitem{Zhao2010}
J.~Zhao, Standard relations of multiple polylogarithm values at roots of
unity, \emph{Documenta Mathematica} \textbf{15} (2010), 1--34.
\href{https://ems.press/journals/dm/articles/8965225}{Publisher record}.
\href{https://arxiv.org/abs/0707.1459}{arXiv:0707.1459}.

\bibitem{DLMF}
NIST Digital Library of Mathematical Functions, \S3.9(ii),
Euler's transformation of series.
\url{https://dlmf.nist.gov/3.9.ii}.
Accessed 9 October 2026. The decrement-difference convention used here
is the negative of DLMF's forward-difference convention.

\bibitem{PSLQ}
H.~R.~P. Ferguson, D.~H. Bailey, and S.~Arno,
Analysis of PSLQ, an integer relation finding algorithm,
\emph{Mathematics of Computation} \textbf{68} (1999), no.~225, 351--369.
\href{https://doi.org/10.1090/S0025-5718-99-00995-3}{doi:10.1090/S0025-5718-99-00995-3}.
\end{thebibliography}
\endgroup
\end{document}
